\begin{tikzpicture}[figure picture,x=10mm,y=10mm,font=\figurefont\fontsize{8}{10}\selectfont]
\path[use as bounding box] (0,0) rectangle (16.4,5.0);
% 顶部：混合测度 ν 在 [0,1] 上
\node[figure heading,anchor=west] at (0.4,4.65) {混合测度 $\nu=\lambda|_{[0,1]}+\delta_0$：同一杆秤上的两笔不同性质的质量};
\draw[FigureGrid,line width=.8pt] (1.0,3.2)--(9.2,3.2);
\fill[FigureBlueFill] (2.2,3.14) rectangle (9.0,3.26);
\draw[FigureBlue,line width=1pt] (2.2,3.14) rectangle (9.0,3.26);
\fill[FigureRed] (2.2,3.2) circle[radius=1.2mm];
\node[below,FigureMuted] at (2.2,3.0) {$0$};
\node[below,FigureMuted] at (9.0,3.0) {$1$};
\node[right,FigureBlue] at (9.05,3.2) {连续铺开的长度};
\node[right,FigureRed] at (2.55,3.45) {点 $0$ 上的整颗质量};
% 分叉箭头到两笔账：蓝色向左下，红色先贴右再折向右下，避免交叉
\draw[FigureBlue,line width=.9pt,-{Stealth[length=2mm]}] (6.4,3.05)--(4.3,1.83);
\draw[FigureRed,line width=.9pt,-{Stealth[length=2mm]}] (2.3,3.0) .. controls (7.0,2.75) and (10.6,2.55) .. (12.1,1.83);
% 左下：绝对连续部分
\fill[FigureBlueFill,rounded corners=2pt] (1.2,0.55) rectangle (7.4,1.8);
\node[FigureBlue,anchor=west] at (1.5,1.5) {第一笔 $\nu_{\mathrm{ac}}=\lambda|_{[0,1]}$};
\node[figure note,anchor=west,align=left] at (1.5,1.0) {能写成密度 $f$、按长度计价；\\参考秤 $\mu$ 称得到};
% 右下：奇异部分
\fill[FigureRedFill,rounded corners=2pt] (9.0,0.55) rectangle (15.2,1.8);
\node[FigureRed,anchor=west] at (9.3,1.5) {第二笔 $\nu_{\mathrm s}=\delta_0$};
\node[figure note,anchor=west,align=left] at (9.3,1.0) {集中在参考零集 $\{0\}$ 上、按点收固定费；\\再大的密度也乘不出这笔};
\end{tikzpicture}
